% ============================================================
\chapter{Actor-Critic Methods}
\label{ch:actor_critic}
% ============================================================

Policy gradient methods (REINFORCE) suffer from high variance; value-based methods (DQN) struggle with continuous action spaces. In the early 2000s, Richard Sutton, David McAllester, Satinder Singh, and Yishay Mansour formalized an idea that reconciles these two families: the \emph{policy gradient theorem} shows that the gradient can be estimated using a separately learned value function. Thus was born the \emph{actor-critic} architecture: an \textbf{actor} that proposes actions and a \textbf{critic} that evaluates their quality. The critic provides a baseline that drastically reduces gradient estimate variance, while the actor retains the ability to represent stochastic policies over continuous spaces. From A2C to PPO to SAC, actor-critic methods dominate contemporary deep reinforcement learning.

\begin{intuition}
Actor-critic methods combine the best of both worlds: an \textbf{actor}
(the policy $\pi_\theta$) and a \textbf{critic} (the value function
$V_\phi$ or $Q_\phi$). The critic provides a low-variance signal to
guide the actor's updates, eliminating the need to wait for complete
episodes. Modern variants (PPO, SAC, TD3) form the backbone of
contemporary deep RL.
\end{intuition}

% ------------------------------------------------------------
\section{Actor-Critic Framework}
\label{sec:ac:framework}
% ------------------------------------------------------------

\begin{definition}[Actor-Critic]
\index{actor-critic}
An \textbf{actor-critic} algorithm maintains two networks:
\begin{itemize}
    \item \textbf{Actor}: $\pi_\theta(a|s)$, parameterized by $\theta$.
    \item \textbf{Critic}: $V_\phi(s)$ or $Q_\phi(s,a)$, parameterized
    by $\phi$.
\end{itemize}
The critic is updated via TD methods and the actor via policy gradient
using the critic as a baseline.
\end{definition}

\begin{keyformulas}[title=\textbf{Basic Actor-Critic Update}]
\textbf{Critic}:
\[
    \phi \leftarrow \phi - \alpha_\phi \nabla_\phi \left(
    R_{t+1} + \gamma V_\phi(S_{t+1}) - V_\phi(S_t) \right)^2
\]
\textbf{Actor}:
\[
    \theta \leftarrow \theta + \alpha_\theta \nabla_\theta
    \log \pi_\theta(A_t|S_t) \,
    \underbrace{\left(R_{t+1} + \gamma V_\phi(S_{t+1})
    - V_\phi(S_t)\right)}_{\hat{A}_t \text{ (estimated advantage)}}
\]
\end{keyformulas}

% ------------------------------------------------------------
\section{Advantage Actor-Critic (A2C)}
\label{sec:a2c}
% ------------------------------------------------------------

\begin{definition}[A2C]
\index{A2C}
Advantage Actor-Critic (A2C) is the synchronous version of A3C
(Mnih et al., 2016). Multiple workers collect trajectories in parallel.
The gradients are averaged and applied synchronously:
\[
    \nabla_\theta J \approx \frac{1}{K} \sum_{k=1}^{K}
    \sum_{t} \nabla_\theta \log \pi_\theta(a_t^{(k)} | s_t^{(k)}) \,
    \hat{A}_t^{(k)}
\]
where $K$ is the number of parallel workers.
\end{definition}

\begin{definition}[A3C]
\index{A3C}
Asynchronous Advantage Actor-Critic (A3C) uses multiple workers updating
a shared parameter server \emph{asynchronously}. Each worker computes
gradients on its own trajectory and applies them to the global parameters
without synchronization.
\end{definition}

\begin{remark}
A2C typically matches or outperforms A3C in practice because synchronous
updates produce less noisy gradients. A3C's main advantage was
computational (avoiding GPU memory bottlenecks), which is less relevant
with modern hardware.
\end{remark}

% ------------------------------------------------------------
\section{Proximal Policy Optimization (PPO)}
\label{sec:ppo}
% ------------------------------------------------------------

\begin{definition}[PPO-Clip]
\index{PPO}
Proximal Policy Optimization (Schulman et al., 2017) uses a clipped
surrogate objective to prevent excessively large policy updates:
\[
    L^{\text{CLIP}}(\theta) = \E_t \left[
    \min\left(r_t(\theta) \hat{A}_t, \;
    \text{clip}(r_t(\theta), 1-\epsilon, 1+\epsilon) \hat{A}_t
    \right) \right]
\]
where $r_t(\theta) = \frac{\pi_\theta(a_t|s_t)}{\pi_{\theta_{\text{old}}}(a_t|s_t)}$
is the probability ratio and $\epsilon \approx 0.2$.
\end{definition}

\begin{keyformulas}[title=\textbf{PPO Full Objective}]
\[
    L(\theta) = L^{\text{CLIP}}(\theta)
    - c_1 \, L^{\text{VF}}(\phi)
    + c_2 \, H[\pi_\theta]
\]
where $L^{\text{VF}} = (V_\phi(s_t) - G_t)^2$ is the value loss,
$H[\pi_\theta]$ is an entropy bonus for exploration, and $c_1, c_2$
are coefficients.
\end{keyformulas}

\begin{algorithme}[title=\textbf{PPO Algorithm}]
\begin{enumerate}
    \item For each iteration:
    \begin{enumerate}
        \item Collect $T$ timesteps of data using $\pi_{\theta_{\text{old}}}$
              with $K$ parallel actors
        \item Compute advantages $\hat{A}_t$ using GAE
        \item For $E$ epochs over the collected data:
        \begin{enumerate}
            \item Sample mini-batch from the collected data
            \item Compute $r_t(\theta)$ and the clipped objective
            \item Update $\theta$ by gradient ascent on $L(\theta)$
        \end{enumerate}
        \item $\theta_{\text{old}} \leftarrow \theta$
    \end{enumerate}
\end{enumerate}
\end{algorithme}

\begin{intuition}
The clipping mechanism in PPO acts as a ``trust region'' in probability
ratio space. When the advantage is positive, the objective is clipped
if $r_t > 1+\epsilon$ (preventing too large a ratio). When the advantage
is negative, it is clipped if $r_t < 1-\epsilon$. This ensures the new
policy stays close to the old one without requiring second-order
optimization.
\end{intuition}

% ------------------------------------------------------------
\section{Continuous Action Spaces: DDPG}
\label{sec:ddpg}
% ------------------------------------------------------------

\begin{definition}[DDPG]
\index{DDPG}
Deep Deterministic Policy Gradient (Lillicrap et al., 2016) extends DQN
to continuous actions using a deterministic policy $\mu_\theta(s)$:
\begin{itemize}
    \item \textbf{Critic}: $Q_\phi(s,a)$ trained with Bellman error:
    $y = r + \gamma Q_{\phi^-}(s', \mu_{\theta^-}(s'))$
    \item \textbf{Actor}: $\theta$ updated by the deterministic policy
    gradient:
    $\nabla_\theta J = \E_s[\nabla_a Q_\phi(s,a)|_{a=\mu_\theta(s)}
    \nabla_\theta \mu_\theta(s)]$
\end{itemize}
Both use target networks with Polyak averaging:
$\phi^- \leftarrow \tau \phi + (1-\tau)\phi^-$.
\end{definition}

\begin{theorem}[Deterministic Policy Gradient]
\index{deterministic policy gradient}
For a deterministic policy $\mu_\theta$, the policy gradient simplifies to:
\[
    \nabla_\theta J(\theta) = \E_{s \sim d^{\mu_\theta}}\left[
    \nabla_\theta \mu_\theta(s) \nabla_a Q^{\mu_\theta}(s,a)\big|_{a=\mu_\theta(s)}
    \right].
\]
This avoids integration over the action space, making it efficient for
high-dimensional continuous actions.
\end{theorem}

% ------------------------------------------------------------
\section{Twin Delayed DDPG (TD3)}
\label{sec:td3}
% ------------------------------------------------------------

\begin{definition}[TD3]
\index{TD3}
TD3 (Fujimoto et al., 2018) addresses DDPG's overestimation with three
techniques:
\begin{enumerate}
    \item \textbf{Twin critics}: two independent Q-networks $Q_{\phi_1},
    Q_{\phi_2}$; the target uses the minimum:
    $y = r + \gamma \min_{i=1,2} Q_{\phi_i^-}(s', \tilde{a}')$
    \item \textbf{Delayed updates}: the actor and target networks are
    updated every $d$ steps (typically $d=2$), while the critics update
    every step.
    \item \textbf{Target policy smoothing}: noise is added to the target
    action: $\tilde{a}' = \mu_{\theta^-}(s') + \text{clip}(\epsilon, -c, c)$,
    $\epsilon \sim \mathcal{N}(0, \sigma^2)$.
\end{enumerate}
\end{definition}

% ------------------------------------------------------------
\section{Soft Actor-Critic (SAC)}
\label{sec:sac}
% ------------------------------------------------------------

\begin{definition}[Maximum Entropy RL]
\index{maximum entropy RL}
The maximum entropy objective augments the standard return with an
entropy bonus:
\[
    J(\theta) = \E_{\pi_\theta}\left[\sum_{t=0}^{\infty} \gamma^t
    \left(R_{t+1} + \alpha H[\pi_\theta(\cdot | S_t)]\right)\right]
\]
where $\alpha > 0$ is the temperature parameter controlling the
exploration-exploitation trade-off.
\end{definition}

\begin{definition}[Soft Actor-Critic]
\index{SAC}
SAC (Haarnoja et al., 2018) learns a stochastic policy by optimizing
the maximum entropy objective:
\begin{itemize}
    \item \textbf{Soft Q-function}:
    $Q(s,a) = r + \gamma \E_{s'}[V(s')]$ where
    $V(s) = \E_{a \sim \pi}[Q(s,a) - \alpha \log \pi(a|s)]$
    \item \textbf{Policy}: $\pi_\theta$ is updated to minimize:
    $\E_{a \sim \pi_\theta}[\alpha \log \pi_\theta(a|s) - Q_\phi(s,a)]$
    \item \textbf{Temperature}: $\alpha$ is automatically tuned to
    maintain a target entropy $\bar{H}$.
\end{itemize}
\end{definition}

\begin{keyformulas}[title=\textbf{SAC Soft Bellman Equation}]
\[
    Q(s,a) = R(s,a) + \gamma \E_{s' \sim P}\left[
    \E_{a' \sim \pi}\left[Q(s', a') - \alpha \log \pi(a'|s')\right]
    \right]
\]
\end{keyformulas}

\begin{proposition}[SAC Convergence]
In the tabular setting, repeated soft policy evaluation and soft policy
improvement converge to the optimal policy of the maximum entropy
objective. The entropy regularization ensures sufficient exploration
and makes the optimization landscape smoother.
\end{proposition}

% ------------------------------------------------------------
\section{Entropy Regularization}
\label{sec:entropy}
% ------------------------------------------------------------

\begin{definition}[Entropy Bonus]
\index{entropy regularization}
Adding an entropy term $H[\pi_\theta(\cdot|s)] = -\sum_a \pi_\theta(a|s)
\log \pi_\theta(a|s)$ to the objective encourages exploration by
penalizing peaked distributions. The coefficient $\alpha$ balances
reward maximization with policy stochasticity.
\end{definition}

\begin{remark}
Entropy regularization has several benefits:
\begin{itemize}
    \item Prevents premature convergence to deterministic policies.
    \item Improves robustness to perturbations.
    \item Enables multi-modal policies that capture multiple near-optimal
    behaviors.
    \item In SAC, automatic temperature tuning adapts $\alpha$ throughout
    training.
\end{itemize}
\end{remark}

% ------------------------------------------------------------
\section{Python Implementation}
\label{sec:ac:code}
% ------------------------------------------------------------

\begin{codeblock}[title=\textbf{PPO (Simplified)}]
\begin{minted}{python}
import torch
import torch.nn as nn
import numpy as np

def ppo_update(policy, value_fn, optimizer,
               states, actions, old_log_probs,
               returns, advantages,
               clip_eps=0.2, epochs=4, batch_size=64):
    states = torch.FloatTensor(np.array(states))
    actions = torch.LongTensor(actions)
    old_lp = torch.FloatTensor(old_log_probs)
    ret = torch.FloatTensor(returns)
    adv = torch.FloatTensor(advantages)
    adv = (adv - adv.mean()) / (adv.std() + 1e-8)

    for _ in range(epochs):
        idx = np.random.permutation(len(states))
        for start in range(0, len(states), batch_size):
            mb = idx[start:start+batch_size]
            probs = policy(states[mb])
            dist = torch.distributions.Categorical(probs)
            new_lp = dist.log_prob(actions[mb])
            ratio = torch.exp(new_lp - old_lp[mb])

            surr1 = ratio * adv[mb]
            surr2 = torch.clamp(ratio, 1-clip_eps,
                                1+clip_eps) * adv[mb]
            policy_loss = -torch.min(surr1, surr2).mean()
            value_loss = nn.MSELoss()(value_fn(states[mb]),
                                      ret[mb])
            entropy = dist.entropy().mean()
            loss = policy_loss + 0.5*value_loss - 0.01*entropy

            optimizer.zero_grad()
            loss.backward()
            nn.utils.clip_grad_norm_(
                list(policy.parameters())
                + list(value_fn.parameters()), 0.5)
            optimizer.step()
\end{minted}
\end{codeblock}

% ------------------------------------------------------------
\section{Exercises}
\label{sec:ac:exercises}
% ------------------------------------------------------------

\begin{exercise}[A2C Implementation]
Implement A2C with GAE on \texttt{CartPole-v1}. Compare different
numbers of parallel workers $K \in \{1, 4, 8, 16\}$ and measure
wall-clock time to convergence.
\end{exercise}

\begin{exercise}[PPO Clipping Analysis]
Visualize the PPO clipped objective $L^{\text{CLIP}}$ as a function of
$r_t$ for $\hat{A}_t > 0$ and $\hat{A}_t < 0$. Explain why the clipping
prevents destructively large updates.
\end{exercise}

\begin{exercise}[SAC on Continuous Control]
Implement SAC for \texttt{Pendulum-v1}. Compare with DDPG and TD3.
Plot Q-value estimates and verify that SAC avoids overestimation.
\end{exercise}

\begin{exercise}[Entropy Temperature Tuning]
Implement automatic entropy temperature tuning for SAC. Compare with
fixed $\alpha \in \{0.01, 0.1, 0.5\}$ and explain the effect on
exploration behavior.
\end{exercise}
